\documentclass[11pt,a4paper]{article}

%% PACHETE MATEMATICE
\usepackage{amsmath,amssymb,amsfonts}
\usepackage{amsthm}
\usepackage{mathpazo}
\usepackage{eulervm}
\usepackage{enumerate}
\usepackage{color}
\usepackage{graphicx}
\usepackage[all]{xy}
\usepackage{xcolor}
\usepackage{framed}
\usepackage[unicode]{hyperref}
\setcounter{MaxMatrixCols}{20}
\usepackage{multicol}
%\usepackage[romanian]{babel}


%% ASPECT
\linespread{1.05}
\usepackage[T1]{fontenc}
\usepackage[utf8x]{inputenc}
\usepackage[margin=2cm]{geometry}
% \usepackage[color]{showkeys}
\usepackage{anyfontsize}
\usepackage{titlesec}
\titleformat*{\section}{\large\bfseries}

\usepackage{fancyhdr}
\pagestyle{fancy}
\rhead{\small \color{gray} Notițe de Adrian Manea}
\lhead{\small \color{gray} ETTI-AN1, 2017---2018, C. Ghiu}


\usepackage[autostyle = false, style = german]{csquotes}
\usepackage[romanian]{babel}
\newcommand\qq{\enquote}

%% COMENZILE MELE
\renewcommand*{\proofname}{Dem.:}
\newcommand\bb{\mathbb}
\newcommand\dr{\mathrm}
\newcommand\kal{\mathcal}
\newcommand\fk{\mathfrak}
\newcommand\op{\oplus}
\newcommand\ot{\otimes}
\newcommand\ds{\displaystyle}
\newcommand\id{\indent}
\newcommand\nid{\noindent}
\newcommand\seq{\subseteq}
\newcommand\sneq{\subsetneq}
\newcommand\speq{\supseteq}
\newcommand\spneq{\supsetneq}
\newcommand\rar{\rightarrow}
\newcommand\Rar{\Rightarrow}
\newcommand\xrar{\xrightarrow}
\newcommand\lar{\leftarrow}
\newcommand\Lar{\Leftarrow}
\newcommand\xlar{\xleftarrow}
\newcommand\lrar{\leftrightarrow}
\newcommand\Lrar{\Leftrightarrow}
\newcommand\ideal{\trianglelefteq}
\newcommand\ai{\text{ a.\^{i}. }}
\newcommand\sk{(\textit{Schi\c{t}\u{a}}) }
\newcommand\rhup{\rightharpoonup}
\newcommand\rhdn{\rightharpoondown}
\newcommand\lhup{\leftharpoonup}
\newcommand\lhdn{\leftharpoondown}
\newcommand\vid{\emptyset}
\newcommand\wh{\widehat}
\newcommand\ol{\overline}
\newcommand\NN{\mathbb{N}}
\newcommand\RR{\mathbb{R}}
\newcommand\ZZ{\mathbb{Z}}
\newcommand\QQ{\mathbb{Q}}
\newcommand\CC{\mathbb{C}}

%% STILURILE MELE
\newtheoremstyle{dotless}{}{}{\itshape}{}{\bfseries}{:}{ }{}

\theoremstyle{dotless}
\newtheorem{theorem}{Teorem\u{a}}[section]
\newtheorem{proposition}{Propozi\c{t}ie}[section]
\newtheorem{lemma}{Lem\u{a}}[section]
\newtheorem{corollary}{Corolar}[section]
\newtheorem{exercise}{Exerci\c{t}iu}

\newtheoremstyle{dotlessdr}{}{}{}{}{\bfseries}{:}{ }{}
\theoremstyle{dotlessdr}
\newtheorem{example}{Exemplu}[section]
\newtheorem{definition}{Defini\c{t}ie}[section]
\newtheorem{remark}{Observa\c{t}ie}[section]




\begin{document}
% \definecolor{labelkey}{RGB}{222,0,0}

\begin{center}
  {\large\textbf{Seminar 3 \\ Serii de funcții. Serii de puteri \\ Exerciții și exemple}}
\end{center}

\section{Serii de numere}
\label{sec:s-num}

1. Studiați convergența seriilor cu termenul general:
\begin{enumerate}[(a)]
\item $x_n = \sqrt{n^4 + 2n + 1} - n^2$; \hfill (comparație la limită)
\item $x_n = \frac{1}{1 + \sqrt{2} + \sqrt[3]{3} + \dots + \sqrt[n]{n}}$; \hfill (comparație)
\item $x_n = (-1)^n \frac{\log_a n}{n}, a > 1$; \hfill (Leibniz)
\item $x_n = \frac{\sin n \cdot \sin n^2}{n^2}$. \hfill (Abel)
\end{enumerate}

2. Arătați că seria $\sum (-1)^n \cdot \frac{1}{n}$ este convergentă, dar
nu este absolut convergentă. (Leibniz)

%%%%%%%%%%%%%%%%%%%%%%%%%%%%%%%%%%%%%%%%%%%%%%%%%%%%%%%%%%%%%%%%%%%%%%

\section{Șiruri de funcții}
\label{sec:s-func}

3. Studiați (cu definiția) convergența punctuală și uniformă pentru șirurile de funcții:
\begin{enumerate}[(a)]
\item $f_n : (0, 1) \to \RR, f_n(x) = \frac{1}{nx + 1}, n \geq 0$;
\item $f_n : [0, 1] \to \RR, f_n(x) = x^n - x^{2n}, n \geq 0$;
\item $f_n : \RR \to \RR, f_n(x) = \sqrt{x^2 + \frac{1}{n^2}}, n > 0$;
\item $f_n : [-1, 1] \to \RR, f_n(x) = \frac{x}{1 + nx^2}$.
\end{enumerate}

%%%%%%%%%%%%%%%%%%%%%%%%%%%%%%%%%%%%%%%%%%%%%%%%%%%%%%%%%%%%%%%%%%%%%%

\section{Serii de funcții}
\label{sec:serii-f}

4. Studiați convergența seriilor de funcții și decideți dacă se pot deriva
termen cu termen:
\begin{enumerate}[(a)]
\item $\sum n^{-x}, x \in \RR$; (Weierstrass, comparație cu armonică)
\item $\sum \frac{\sin nx}{2^n}, x \in \RR$;
\item $\frac{\sin nx}{n(n+1)}$.
\end{enumerate}

\emph{Soluții:}
(a) Seria converge punctual dacă și numai dacă $x > 1$, pentru că poate fi
comparată cu seria armonică generalizată.

Fie, deci $r > 1$. Pentru orice $x \geq r$, avem $\frac{1}{n^x} \leq \frac{1}{n^r}$.
Seria $\sum \frac{1}{n^r}$ este convergentă și rezultă din criteriul lui
Weierstrass că seria converge uniform pe intervalul $[r, \infty)$.

Pentru seria derivatelor: $\sum (n^{-x})' = -\sum n^{-x} \ln n$, găsim că ea
converge uniform pe orice interval $[r, \infty)$. Într-adevăr, putem aplica
criteriul lui Weierstrass:
\[
  \sum_{n \geq 1} n^{-x} \ln n \leq \sum_{n \geq 1} n^{-r} \ln n, \forall x \geq r.
\]
A doua serie, privită ca o serie numerică, este convergentă.

(b) Deoarece avem $|f_n(x)| \leq \frac{1}{2^n}, \forall x \in \RR$, iar seria
obținută este seria geometrică, cu rație subunitară, deci convergentă, deducem
din criteriul lui Weierstrass că seria este uniform convergentă pe $\RR$.

Seria derivatelor:
\[
  \sum_{n \geq 1} \Big( \frac{\sin nx}{2^n} \Big) = %
  \sum_{n \geq 1} \frac{n \cos nx}{2^n},
\]
care converge uniform pe $\RR$, din criteriul lui Weierstrass. Așadar, seria se
poate deriva termen cu termen.

(c) Putem folosi din nou criteriul lui Weierstrass:
\[
  \sum_{n \geq 1} \frac{\sin nx}{n(n + 1)} \leq %
  \sum_{n \geq 1} \frac{1}{n(n+1)},
\]
iar a doua serie este convergentă. Așadar, seria inițială este uniform convergentă.

Seria derivatelor $\sum \frac{\cos nx}{n + 1}$ nu converge punctual (comparație
cu seria armonică, de exemplu), deci nu se poate deriva termen cu termen.

\vspace{1cm}

5. Să se arate că șirul de funcții $(f_n)_n$, cu:
\[
  f_n : [0, 1] \to \RR, \quad f_n(x) = nxe^{-nx^2}
\]
converge, dar avem:
\[
  \lim_{n \to \infty} \int_0^1 f_n(x) dx \neq \int_0^1 \lim_{n \to \infty} f_n(x) dx.
\]

\emph{Soluție:} Avem $\ds\lim_{n \to \infty}f_n(x) = 0$, deci $f(x) = 0$, pentru
$x \in [0, 1]$. Calculăm integrala:
\[
  \int_0^1 f_n(x) dx = \int_0^1 nxe^{-nx^2} dx = -\frac{1}{2} e^{-nx^2}\Big|_0^1 = %
  \frac{1}{2} - \frac{1}{2} e^{-n}.
\]

Obținem, atunci:
\[
  \lim_{n \to \infty} \int_0^1 f_n(x) dx = \frac{1}{2},
\]
dar $\ds\int_0^1 \lim_{n \to \infty} f_n(x) dx = 0$.

Rezultatul se explică prin faptul că șirul nu este uniform convergent. Într-adevăr,
dacă $x_n = \frac{1}{n} \in [0, 1]$, rezultă că $f_n(x_n) = e^{-\frac{1}{n}} \to 1$,
deci definiția convergenței uniforme nu este verificată.

%%%%%%%%%%%%%%%%%%%%%%%%%%%%%%%%%%%%%%%%%%%%%%%%%%%%%%%%%%%%%%%%%%%%%%

\section{Serii de puteri. Serii Taylor}
\label{sec:taylor}

6. Să se dezvolte în seria Maclaurin următoarele funcții, precizînd și
domeniul de convergență:
\begin{enumerate}[(a)]
\item $f(x) = e^x$;
\item $f(x) = \sin x$.
\end{enumerate}

\emph{Soluție:} (a) Calculăm derivatele funcției și obținem că $(e^x)^{(n)} = e^x$,
pentru orice $n$. Atunci, din formula lui Maclaurin, găsim:
\[
  e^x = \sum_{n \geq 0} \frac{1}{n!} \cdot x^n.
\]
Pentru a determina domeniul de convergență, folosim criteriul raportului:
\[
  \lim_{n \to \infty} \Big| \frac{x^{n+1}}{(n+1)!} \cdot \frac{n!}{x^n} \Big| = %
  \lim_{n \to \infty} \Big| \frac{x}{n+1} \Big| = 0 < 1,
\]
deci seria este convergentă pe $\RR$.

Restul de ordin $n$ se obține:
\[
  R_n(x) = \frac{x^{n + 1}}{(n + 1)!} \cdot e^{\xi}, \quad \xi \in (0, x) \text{ sau } %
  \xi \in (x, 0).
\]

(b) Calculăm derivata de ordin $n$ a funcției sinus și observăm:
\[
  (\sin x)^{(n)} = \sin \Big( x + n \frac{\pi}{2} \Big), \forall n \in \NN.
\]
Deducem că derivatele de ordin par sînt nule, iar cele de ordin impar sînt $(-1)^n$.
Rezultă:
\[
  \sin x = \sum_{n \geq 0} \frac{(-1)^n}{(2n + 1)!} \cdot x^{2n + 1}.
\]
Pentru determinarea domeniului de convergență, putem folosi din nou criteriul
raportului și obținem:
\[
  \lim_{n \to \infty} \Big| \frac{x^{2n + 3}}{(2n + 3)!} \cdot \frac{(2n + 1)!}{x^{2n + 1}} \Big| = %
  \lim_{n \to \infty} \frac{x^2}{(2n + 2)(2n + 3)} = 0 < 1,
\]
deci seria este convergentă pe $\RR$.

\vspace{1cm}

7. Să se dezvolte în serie de puteri ale lui $x - 1$ funcția:
\[
  f : \RR - \{ 0 \} \to \RR, \quad f(x) = \frac{1}{x}.
\]

\emph{Soluție:} Derivata de ordin $n$ este:
\[
  f^{(n)}(x) = \frac{(-1)^n n!}{x^{n+1}} \Rightarrow f^{(n)}(1) = (-1)^n n!, \forall n \in \NN.
\]
Rezultă:
\[
  \frac{1}{x} = \sum_{n \geq 0} (-1)^n (x - 1)^n, \forall x \in (0, 2).
\]

\vspace{1cm}

8. Să se afle raza de convergență și mulțimea de convergență pentru seriile de puteri:
\begin{enumerate}[(a)]
\item $\ds\sum_{n \geq 0} x^n$;
\item $\ds\sum_{n \geq 1} \frac{n^n x^n}{n!}$;
\end{enumerate}

\emph{Soluție:} (a) Fie $R$ raza de convergență. Din definiție, avem:
\[
  R = \lim_{n \to \infty} \frac{|a_n|}{|a_{n + 1}|} = 1.
\]
Deducem că seria este absolut convergentă pe $(-1, 1)$ și divergentă în rest.

Evident, seria este uniform convergentă pe orice interval închis $|x| \leq r < 1$, iar
pentru $x = \pm 1$, obținem o serie divergentă.

(b) Din nou, calculăm raza de convergență:
\[
  R = \lim_{n \to \infty} \frac{|a_n|}{|a_{n + 1}|} = %
  \lim_{n \to \infty} \Big( \frac{n}{n + 1} \Big)^n = \frac{1}{e}.
\]
Rezultă că seria este absolut convergentă pe $\Big( -\frac{1}{e}, \frac{1}{e} \Big)$
și divergentă în rest. Seria este uniform convergentă pe $|x| \leq r < e$, iar
pentru $x = \pm \frac{1}{e}$, se obține o serie divergentă, rezultat care
se poate demonstra cu criteriul raportului.

\vspace{1cm}

9. Studiați convergența seriilor de funcții:
\begin{enumerate}[(a)]
\item $\ds\sum_{n \geq 0} \ln^n x, x > 0$;
\item $\ds\sum_{n \geq 1} \Big( 1 + \frac{1}{n} \Big)^{-n^2} e^{-nx}$;
\item $\ds\sum_{n \geq 1} \frac{1}{n! x^n}$.
\end{enumerate}

\emph{Soluție:} Exercițiile se rezolvă cu o schimbare de variabilă, care
transformă seriile de funcții în serii de puteri:
\begin{enumerate}[(a)]
\item $t = \ln x$;
\item $t = e^{-x}$;
\item $t = x^{-1}$.
\end{enumerate}

\vspace{1cm}

10. Să se calculeze sumele seriilor de funcții:
\begin{enumerate}[(a)]
\item $\ds\sum_{n \geq 0} x^{2n}$;
\item $\ds\sum_{n \geq 0} (-1)^n \frac{x^{2n + 1}}{2n + 1}$;
\item $\ds\sum_{n \geq 0} (n + 1)x^n$.
\end{enumerate}

\emph{Soluție:} (a) Putem folosi direct formula pentru seria geometrică:
\[
  \sum_{n \geq 0} x^{2n} = \sum_{n \geq 0} (x^2)^n = \frac{1}{1 - x^2}, %
  \forall x \in (-1, 1).
\]

(b) Raza de convergență se poate calcula și obținem $R = 1$. Fie $S(x)$ suma seriei
date. Seria derivatelor are aceeași rază de convergență cu seria inițială, iar
suma ei este egală cu derivata sumei (seriile de puteri se pot deriva mereu termen
cu termen):
\[
  S'(x) = \sum_{n \geq 0} (-1)^n (x^2)^n = \frac{1}{1 + x^2}, \forall x \in (-1, 1).
\]
Rezultă $S(x) = \arctan x + C$ și, deoarece $S(0) = 0$, avem $C = 0$.

(c) Procedăm similar, acum cu integrare termen cu termen. Raza de convergență
este egală cu 1. Avem:
\[
  \int_0^x S(x) dx = \sum_{n \geq 0} x^{n + 1} = \frac{x}{1 - x}, \forall x \in (-1, 1).
\]
Rezultă $S(x) = \frac{1}{(1 - x)^2}$.

%%%%%%%%%%%%%%%%%%%%%%%%%%%%%%%%%%%%%%%%%%%%%%%%%%%%%%%%%%%%%%%%%%%%%%

\section{Exerciții propuse}
\label{sec:ex}

1. Să se studieze convergența punctuală și uniformă a șirurilor de funcții:
\begin{enumerate}[(a)]
\item $f_n : \RR \to \RR, f_n(x) = \frac{nx^2}{1 + nx^2}$;
\item $f_n : [0, 1] \to \RR, f_n(x) = \frac{x^n}{1 + x^{2n}}$;
\item $f_n : \RR_+ \to \RR, f_n(x) = e^{-nx} \sin nx$;
\end{enumerate}

\vspace{1cm}

2. Să se afle mulțimea de convergență a seriei de funcții:
\[
  \sum_{n \geq 1} (-1)^n \frac{1}{n^x}.
\]

\vspace{1cm}

3. Să se studieze convergența seriilor:
\begin{enumerate}[(a)]
\item $\ds\sum_{n \geq 1} \frac{1}{x^2 + n^2}$;
\item $\ds\sum_{n \geq 1} \frac{\sin nx}{\sqrt[3]{n^4 + x^4}}$;
\item $\ds\sum_{n \geq 1} \frac{1}{x^2 + 2^n}, x \in \RR$.
\end{enumerate}

\vspace{1cm}

4. Să se determine mulțimea de convergență și suma seriilor:
\begin{enumerate}[(a)]
\item $\ds \sum_{n \geq 1} (-1)^{n + 1} \frac{x^n}{n}$;
\item $\ds \sum_{n \geq 0} (n + 1)(n + 2)(n + 3) x^n$;
\item $\ds \sum_{n \geq 1} n^3 x^n$;
\end{enumerate}

\vspace{1cm}

5. Găsiți mulțimea de convergență pentru seriile de puteri:
\begin{enumerate}[(a)]
\item $\ds\sum_{n \geq 1} \frac{1}{n! x^n}$;
\item $\ds\sum_{n \geq 1} \frac{x^n}{2^n + 3^n}$;
\item $\ds\sum_{n \geq 1} \sqrt[n]{n!} x^n$;
\item $\ds\sum_{n \geq 1} \frac{1 \cdot 5 \cdot 9 \cdots (4n - 3)}{ %
    3 \cdot 7 \cdot 11 \cdots (4n - 1)} x^n$;
\end{enumerate}

\vspace{1cm}

6. Să se dezvolte în serie Maclaurin următoarele funcții, precizînd domeniul
de convergență:
\begin{enumerate}[(a)]
\item $f(x) = \cos x$;
\item $f(x) = (1 + x)^a, a \in \RR$;
\item $f(x) = \frac{1}{1 + x}$;
\item $f(x) = \sqrt{1 + x}$;
\item $f(x) = \ln ( 1 + x)$;
\item $f(x) = \arctan x$.
\end{enumerate}

\vspace{1cm}

7. Să se dezvolte în serie de puteri ale lui $x + 4$ funcția:
\[
  f: \RR - \{ -2, -1 \} \to \RR, \quad %
  f(x) = \frac{1}{x^2 + 3x + 2}.
\]



\end{document}